\section{变分推断把后验近似变成优化}
\label{sec:13-Machine-Learning/08-Sampling-and-Inference/03}

推荐系统的团队要给两千万用户各拟合一个隐向量的后验，模型每晚重训一次，第二天早上八点必须上线。上一节那条链在这里根本跑不起来：一次密度评估要扫过全部交互记录，走一万步就是上亿次扫描，何况还得跑多条链做诊断。就算算力管够，后验的形状也未必配合——两个被低密度谷隔开的峰，链停在其中一个上，trace 平得像直线，而你算出的后验均值系统性地偏在一边。

变分推断换掉了整个问法。它不打算从后验里抽样，而是问：在一族容易处理的分布 \(\mathcal Q\) 里，哪个成员 \(q_\phi(z)\) 离真实后验 \(p(z\given x)\) 最近。推断于是从采样问题变成优化问题，而优化是有收敛曲线、有提前终止、有 minibatch 的——所有工程手段都能用上。这笔交易让出的东西从第一步就写进了公式，只是不容易被注意到。

\subsection{ELBO 把不可算的证据换成可算的下界}

从一个恒等式出发。对任意分布 \(q_\phi(z)\)，

\[
\log p(x)
=\E_{q_\phi}\!\left[\log\frac{p(x,z)}{q_\phi(z)}\right]
+\E_{q_\phi}\!\left[\log\frac{q_\phi(z)}{p(z\given x)}\right].
\]

第一项记作

\[
\mathcal L(\phi)=\E_{q_\phi}[\log p(x,z)-\log q_\phi(z)],
\]

称为证据下界；第二项恰是 \(\KL(q_\phi\|p(\cdot\given x))\)。于是

\[
\log p(x)=\mathcal L(\phi)+\KL\bigl(q_\phi(z)\,\|\,p(z\given x)\bigr).
\]

左端与 \(\phi\) 无关，KL 非负，两条一合，最大化 ELBO 精确等价于最小化 \(\KL(q_\phi\|p)\)，而 ELBO 永远不超过证据，相等当且仅当 \(q_\phi\) 就是后验。

这个恒等式真正解决的是可计算性。\(\mathcal L(\phi)\) 只用到联合分布 \(p(x,z)\) 与变分分布本身，两者都能逐点求值；被绕开的仍然是那个归一化常数 \(p(x)\)。上一节靠接受率把它约掉，这里靠下界把它挡在等式另一侧——同一个障碍，两种回避方式。

失掉的东西也写在同一行里。即使优化做到极致，你拿到的是 \(\mathcal Q\) 中最好的那个成员，不是后验。变分族在你写下它的那一刻，就决定了哪些后验形状永远表达不出来。

\subsection{Mean field 用独立性换来闭式更新}

最常用的族假设隐变量分组之间彼此独立：

\[
q(z)=\prod_{j=1}^m q_j(z_j).
\]

这一步把真实后验里 \(z_1\) 与 \(z_2\) 的相关性一笔勾销，换来的是逐因子优化有闭式解。把 ELBO 中含 \(q_j\) 的部分拎出来，

\[
\mathcal L=\E_{q_j}\bigl[\E_{q_{-j}}\log p(x,z)\bigr]-\E_{q_j}[\log q_j(z_j)]+\text{与 }q_j\text{ 无关的项},
\]

这正是 \(-\KL(q_j\|\tilde p_j)\) 加常数，其中 \(\tilde p_j(z_j)\propto\exp\bigl(\E_{q_{-j}}[\log p(x,z)]\bigr)\)。KL 在两者相等时取零，于是最优因子为

\[
\log q_j^{*}(z_j)=\E_{q_{-j}}[\log p(x,z)]+\text{常数},
\]

常数由归一化定出。条件对账只有两条：期望 \(\E_{q_{-j}}[\log p(x,z)]\) 要存在且可算，指数化之后要可归一化。指数族配共轭先验时两条自动满足，\(q_j^{*}\) 与先验同族，逐因子轮转更新即坐标上升变分推断，ELBO 单调不降。

这套更新与 \hyperref[sec:13-Machine-Learning/06-Probabilistic-Models/04]{``EM 用下界交替完成隐变量学习''} 的 E 步长得很像，角色却不同。EM 的 E 步用的是当前参数下的精确后验，M 步把参数当点估计推进；变分推断可以给参数也保留一个分布，只是这个分布被关在独立性假设里。

\subsection{重参数化让梯度穿过采样这一步}

期望没有解析形式时，ELBO 的梯度得用 Monte Carlo 估。麻烦在于 \(\nabla_\phi\E_{q_\phi}[f(z)]\) 不能直接把梯度搬进期望号——积分的测度本身依赖 \(\phi\)。

重参数化把随机性与参数拆开。若存在变换 \(z=g_\phi(\varepsilon)\)，\(\varepsilon\sim p(\varepsilon)\)，使 \(z\) 的分布恰为 \(q_\phi\) 且 \(\varepsilon\) 的分布与 \(\phi\) 无关，则

\[
\nabla_\phi\E_{q_\phi}[f(z)]=\E_{p(\varepsilon)}\bigl[\nabla_\phi f(g_\phi(\varepsilon))\bigr],
\]

期望的测度不再动，梯度顺着确定性变换流回参数。Gaussian 族是最干净的例子：\(z=\mu+\sigma\odot\varepsilon\)，\(\varepsilon\sim\mathcal N(0,I)\)，参数全部搬进了变换里，采样只发生在与 \(\phi\) 无关的标准正态上。

到这里，误差已经分成四层，混起来看会导出错误的补救动作。优化误差是还没找到 \(\mathcal Q\) 里 ELBO 最大的那个成员，加迭代、换优化器、调步长有用；Monte Carlo 梯度误差来自有限样本，加大批量或用方差更低的估计量有用；近似误差是 \(\mathcal Q\) 压根表达不出真实后验的形状，只有换族才有用；模型误差是 \(p(x,z)\) 与数据生成机制的差距，换推断方法一点用都没有。看到不确定性偏小就去加迭代次数，通常是把第三层错误当成第一层在治。

\subsection{反向 KL 的模式寻找让方差被系统性低估}

变分推断最小化的是 \(\KL(q\|p)\)，期望取在 \(q\) 上。这个方向有很具体的脾气。若 \(p\) 在某处几乎为零而 \(q\) 在那里放了质量，\(\log(q/p)\) 爆掉，惩罚极重；反过来，\(p\) 在某处质量可观而 \(q\) 近乎为零时，被积式 \(q\log(q/p)\) 里的 \(q\) 本身已经很小，惩罚可以忽略。两种错误的代价不对称，\(q\) 于是宁可缩进 \(p\) 的一块高密度区，也不肯把质量摊到 \(p\) 不支持的地方。KL 两个方向的不对称在 \hyperref[sec:07-Information-Theory/03-Divergences/02]{``Gibbs 不等式与 KL 的方向性''} 里有完整讨论，这里看到的是它在推断中的直接后果。

\begin{center}
\begin{tikzpicture}[>=stealth]
\draw[->] (-4.2,0) -- (4.4,0) node[right]{\(z\)};
\draw[very thick] plot[smooth] coordinates
  {(-3.6,0.03) (-3.0,0.14) (-2.6,0.50) (-2.2,1.05) (-2.0,1.25)
   (-1.8,1.05) (-1.4,0.50) (-1.0,0.14) (-0.6,0.05) (0,0.02)
   (0.6,0.05) (1.0,0.14) (1.4,0.50) (1.8,1.05) (2.0,1.25)
   (2.2,1.05) (2.6,0.50) (3.0,0.14) (3.6,0.03)};
\draw[dashed,thick] plot[smooth] coordinates
  {(-3.4,0.02) (-3.0,0.10) (-2.7,0.36) (-2.4,0.80) (-2.0,1.18)
   (-1.6,0.80) (-1.3,0.36) (-1.0,0.10) (-0.6,0.02)};
\draw[dotted,thick] plot[smooth] coordinates
  {(-3.9,0.08) (-3.2,0.20) (-2.6,0.34) (-2.0,0.46) (-1.2,0.56)
   (-0.6,0.61) (0,0.63) (0.6,0.61) (1.2,0.56) (2.0,0.46)
   (2.6,0.34) (3.2,0.20) (3.9,0.08)};
\draw[very thick] (-3.9,2.15) -- (-3.3,2.15);
\node[right] at (-3.2,2.15) {真实后验 \(p\)};
\draw[dashed,thick] (-3.9,1.8) -- (-3.3,1.8);
\node[right] at (-3.2,1.8) {\(\argmin\KL(q\|p)\)：钻进一个模态};
\draw[dotted,thick] (-3.9,1.45) -- (-3.3,1.45);
\node[right] at (-3.2,1.45) {\(\argmin\KL(p\|q)\)：摊平覆盖全部};
\draw[<->,gray] (0.5,-0.35) -- (3.7,-0.35);
\node[gray] at (1.7,-0.72) {被反向 KL 判为不存在的那一半};
\end{tikzpicture}
\end{center}

\emph{同一个后验、同一个单峰变分族，换个 KL 方向就得到宽窄相反的两个答案。}

图里那条虚线就是双峰后验上的典型结局：单峰 Gaussian 贴住左边那个峰，给右峰的质量近乎为零，宽度也只反映左峰内部的散布。你看到的是 ELBO 平稳收敛、变分参数不再变动、梯度范数降到很小——所有指标都健康——而后验方差被低估了不止一个数量级。这不是实现里有 bug，是目标函数选了这个方向就会这样。想要点线那种覆盖行为，得去最小化 \(\KL(p\|q)\)，可那个期望取在 \(p\) 上，正是你采不了样的分布，期望传播与重要性加权的变分方法都在这条路上付学费。

Mean field 从另一个方向削同一件东西。就算每个 \(z_j\) 的边缘位置合理，因子之间的协方差被硬置为零。真实后验里同涨同跌的两个参数，在近似里各走各的，联合决策要用的那个方差——比如两个效应之和的不确定性——可能小掉一个量级。加大族容量能缓解：全协方差 Gaussian、Gaussian 混合、normalizing flow、按模型结构保留部分依赖。扩大族不会让最优 ELBO 变差，但目标非凸，实际找到的解可能更不稳定，训练成本也上去了。

\subsection{ELBO 不能告诉你近似够不够用}

ELBO 上升只说明 \(q\) 比自己之前更接近后验，它不给出还差多远。原因就在那个恒等式里：\(\log p(x)-\mathcal L(\phi)=\KL(q_\phi\|p)\)，而 \(\log p(x)\) 恰好是你算不出的那个量。缺了它，ELBO 的绝对数值不能跨模型、跨实现比较，它是优化器的手柄，不是裁判的记分牌。

能做的检查都是外部的。在能精确算出后验的简化版本上——小数据、共轭模型——直接比边缘分布、相关系数和尾概率；或者跑一条经过充分诊断的长链作参考，比较同样几个量。用多个分散初值重跑优化，若五个起点给出五组结论，问题还停在优化那一层。逐级加大族容量，从 mean field 到全协方差再到混合，看关键结论稳不稳；加一个混合成分就让某个效应反号，说明原来的结论本就靠不住。还有一个便宜的信号：把学到的 \(q\) 当提议分布做重要性采样，看权重的有效样本量，权重塌缩说明 \(q\) 的尾部比后验薄得多，这正是下一节要反复用到的诊断。

判断的落点是用途。下游只要后验预测均值时，低估协方差的影响有限；要拿 \(99\%\) 分位数决定停不停机，同一个近似就不能用。维数高、数据量大、后验大体单峰、精度要求适中，变分推断是划算的；模态多、尾部要准、算力允许，长链加多起点仍然是更可信的选择。两条路都靠不住的时候还有第三种处境：模型本身就没有一个明确的后验可谈，你要的只是估计量在重复抽样下的波动——下一节从这里出发。
